% ===========================================================================
% Part XII - Integrated Mini Project and Appendices
% ===========================================================================

\part{Integrated Mini Project and Appendices}

% ---------------------------------------------------------------------------
\chapter{Integrated Mini Project: Comparative Machine Learning Study}\label{exp:mini}

\quickfacts
  {Scope}{Compare algorithms from different families on identical data, splits and seeds}
  {Protocol}{Fixed random seeds; untouched test sets; two or more metrics per task}
  {Headline result}{Gradient boosting leads tabular regression; ensembles tie on small classification data; density clustering finds real seismic zones without labels}

\section{Objective}
For the final laboratory assessment, compare at least three algorithms from different
families on one real-world dataset, with a controlled, explainable experiment rather than a
search for the largest score.

\section{Suggested Dataset}
The manual's own experiments provide the comparison material: Heart Disease (classification),
California Housing (regression), Mall Customers (clustering) and the USGS earthquake
catalogue (density-based clustering).

\section{Required Libraries}
scikit-learn, xgboost, catboost, numpy, pandas, matplotlib.

\placeholder

\section{Workflow}
\section{Comparative Results}
\section{Discussion and Conclusions}
\section{Viva / Discussion Questions}

% ===========================================================================
% Appendices
% ===========================================================================

\chapter{Appendix A: Quick Reference Formulas}

\begin{longtable}{@{}p{0.30\textwidth}p{0.64\textwidth}@{}}
\toprule
\textbf{Concept} & \textbf{Formula / Definition} \\
\midrule
\endhead
Euclidean distance & $d(x,y) = \sqrt{\sum_j (x_j - y_j)^2}$ \\
K-means objective & $J = \sum_i \lVert x_i - \mu_{c_i} \rVert^2$ \\
Regression MAE & $\mathrm{MAE} = \frac{1}{n}\sum_i |y_i - \hat{y}_i|$ \\
Regression RMSE & $\mathrm{RMSE} = \sqrt{\frac{1}{n}\sum_i (y_i - \hat{y}_i)^2}$ \\
$R^2$ & $R^2 = 1 - \mathrm{SS_{res}} / \mathrm{SS_{tot}}$ \\
Precision & $\mathrm{TP} / (\mathrm{TP} + \mathrm{FP})$ \\
Recall & $\mathrm{TP} / (\mathrm{TP} + \mathrm{FN})$ \\
F1 & $2 \times \mathrm{Precision} \times \mathrm{Recall} / (\mathrm{Precision} + \mathrm{Recall})$ \\
Sigmoid & $\sigma(z) = 1 / (1 + e^{-z})$ \\
SVM margin idea & Maximize margin while penalizing violations via $C$ \\
GRNN kernel & $w_i(x) = \exp\!\left(-\lVert x - x_i \rVert^2 / (2\sigma^2)\right)$ \\
\bottomrule
\end{longtable}

\chapter{Appendix B: Reproducibility Checklist}

\begin{itemize}
  \item Python version recorded.
  \item Package versions recorded.
  \item Random seed recorded where supported (all experiments use seed 42).
  \item Dataset source and version recorded (public datasets committed under \texttt{data/}).
  \item Preprocessing steps recorded (scaling fitted on training data only; haversine radians for geospatial data).
  \item Train/validation/test protocol recorded (stratified, chronological, early-stopping split, 5-fold CV).
  \item Hyperparameters recorded.
  \item Evaluation metrics defined before final testing.
  \item Plots and tables saved with meaningful names.
  \item Code runs from a clean environment or requirements file.
\end{itemize}

\chapter{Appendix C: Suggested Assessment Rubric}

\begin{longtable}{@{}p{0.68\textwidth}r@{}}
\toprule
\textbf{Component} & \textbf{Suggested Weight} \\
\midrule
\endhead
Pre-lab preparation / algorithm understanding & 10\% \\
Correct implementation & 25\% \\
Experimental design and preprocessing & 15\% \\
Evaluation and visualization & 15\% \\
Analysis and interpretation & 20\% \\
Report quality / reproducibility & 10\% \\
Viva / discussion & 5\% \\
\bottomrule
\end{longtable}

\chapter{Appendix D: Student Experiment Record}

\begin{longtable}{@{}p{0.36\textwidth}p{0.58\textwidth}@{}}
\toprule
\textbf{Field} & \textbf{Student Entry} \\
\midrule
\endhead
Name & \\
ID & \\
Section & \\
Experiment No. & \\
Date & \\
Dataset & \\
Algorithm / Version & \\
Key Hyperparameters & \\
Random Seed & \\
Main Result & \\
Observation & \\
Instructor Signature & \\
\bottomrule
\end{longtable}
